\clearpage
\section{Construct Validity Checklist}

For usability, we have reproduced the complete construct validity checklist here, grouped by recommendation. We recommend that checklist users consider each of these questions and whether the answer is adequately addressed in the benchmark and corresponding paper. We anticipate that it may be difficult to adopt every recommendation in every case. For example, computing confidence intervals may be prohibitively expensive. In these cases, considering and discussing the tradeoffs of adopting the recommendations as limitations would enable readers of these papers to better interpret the results.

\textbf{Define the phenomenon}
\begin{itemize}[leftmargin=25pt, label=\scriptsize$\square$]
        %\item The phenomenon being measured is clearly defined.
        \item Provide a precise and operational definition for the phenomenon being measured
        %\item The benchmark captures the entire targeted phenomenon, not a single aspect of it.
        \item Specify the scope of the phenomenon being covered and acknowledge any excluded aspects
        %\item When the benchmark aggregates sub-aspects of the phenomenon, they are described and measured separately.
        \item Identify if the phenomenon has sub-components and ensure they are measured separately
    \end{itemize}

\textbf{Measure only the phenomenon}
\begin{itemize}[leftmargin=25pt, label=\scriptsize$\square$]
    \item Control for unrelated tasks that may affect the results
%\item If the benchmark has a strict output format, model performance is compared with and without format constraints.
\item Assess the impact of format constraints on model performance
%\item If the benchmark uses automated techniques to parse answers, they are tested for bias and accuracy across different models.
\item Validate any automated output parsing techniques for accuracy, consistency and bias

\end{itemize}

\textbf{Construct a representative dataset for the task}
\begin{itemize}[leftmargin=25pt, label=\scriptsize$\square$]
    \item Employ sampling strategies to ensure task items are representative of the overall task space
 %\item Checks are performed to ensure that the task items are high quality and related to the phenomenon (especially if the benchmark is large).
 \item Verify the quality and relevance of all task items, especially for large or automatically generated datasets
 %\item The selection of task items has been tailored for testing LLMs (e.g. including both easy and difficult tasks by human standards, including small syntactic perturbations of task items).
 \item Include task items that test known LLM sensitivities (e.g. input permutations or variations)
\end{itemize}

\textbf{Acknowledge limitations of reusing datasets}
\begin{itemize}[leftmargin=25pt, label=\scriptsize$\square$]
    \item Document whether the benchmark adapts a previous dataset or benchmark
    %\item If so, the authors provide a clear analysis of the previous work describing strengths and limitations of the benchmark.
    \item If so, analyse and report the relevant strengths and limitations of the adapted prior work
    %\item If so, results on the original benchmark are included and compared to.
    \item If so, report and compare performance on the new benchmark against the original
    %\item Any differences from the original dataset are justified and explained in the context of construct validity.
    %\item Articulate the rationale when modifying a reused dataset, identifying improvements to construct validity
    \item Explain modifications to reused datasets and how they improve construct validity
\end{itemize}

\textbf{Prepare for contamination}
\begin{itemize}[leftmargin=25pt, label=\scriptsize$\square$]
    \item Implement tests to detect data contamination and apply them to the benchmark
    %\item Held-out task items are available (e.g. via a private test set, or generating new questions).
    \item Maintain a held-out set of task items to facilitate ongoing, uncontaminated evaluation
    %\item The authors conduct an analysis of data exposure prior to the benchmark creation.
    \item Investigate the potential pre-exposure of benchmark source materials or similar data in common LLM training corpora
\end{itemize}

\textbf{Use statistical methods to compare models}
\begin{itemize}[leftmargin=25pt, label=\scriptsize$\square$]
     \item Report the benchmark's sample size and justify its statistical power
     %\item Uncertainty estimates are provided for the main scores, and are narrow enough to meaningfully compare relevant models.
     \item Report uncertainty estimates for all primary scores to enable robust model comparisons
     %\item If human raters are used, the recruitment accounts for demographic biases which may be relevant to the preferences they report.
     \item If using human raters, describe their demographics and mitigate potential demographic biases in rater recruitment and instructions
    %\item If the benchmark uses subjective measures of performance, the distribution of labels is reported and accounted for in the scoring.
    \item Use metrics that capture the inherent variability of any subjective labels, without relying on single-point aggregation or exact matching.
\end{itemize}

\textbf{Conduct an error analysis}
\begin{itemize}[leftmargin=25pt, label=\scriptsize$\square$]
    \item Conduct a qualitative and quantitative analysis of common failure modes
    %\item There are no patterns of failures which relate to non-targeted phenomena.
    \item Investigate whether failure modes correlate with non-targeted phenomena (confounders) rather than the intended construct
    %\item If so, potential biases in the scoring are discussed.
    \item If so, identify and discuss any potential scoring biases revealed in the error analysis
    %\item Experiments or recommendations are made to improve model scores on the benchmark.
    \item Conduct experiments or propose new directions to improve model scores on the benchmark
\end{itemize}

\textbf{Justify construct validity}
\begin{itemize}[leftmargin=25pt, label=\scriptsize$\square$]
    \item Justify the relevance of the benchmark for the phenomenon with real-world applications
    %\item The authors explain why the task and metric were chosen to measure the target phenomenon.
    \item Provide a clear rationale for the choice of tasks and metrics, connected to the operational definition of the phenomenon
    %\item The benchmark is compared to other benchmarks of similar phenomena, with a discussion of similarities and differences.
    \item Compare similarities and differences between the benchmark and existing evaluations of similar phenomena
    \item Discuss the limitations and design trade-offs of the benchmark concerning construct validity
\end{itemize}
